\chapter{Finite-Galois descent to the original field}
\label{route-ch-galois}

Let \(C/K\) be the smooth proper geometrically irreducible curve fixed
above. Fix the witness \((P,L,\xi_L)\) of
Theorem~\ref{route-thm-descended-universal-class}, with \(L/K\) finite
Galois, and the natural equivalence of
Theorem~\ref{route-thm-finite-stage-yoneda}:
\[
  r_L:h_{J_L}\xrightarrow{\sim}\operatorname{Pic}^0_{C_L},
  \qquad J_L=P.\mathrm{gluedOver}.
\]
The group structure, affine neighbourhoods of Galois orbits, and descent
below all concern this same \(L\)-scheme \(J_L\).

\section{Quasi-projectivity and finite affine neighbourhoods}

\begin{lemma}[The finite-stage carrier is an algebraic group]
  \label{route-lem-exact-carrier-group}
  \notready
  \uses{route-thm-finite-stage-yoneda, route-thm-glue-package}

  The equivalence \(r_L\) gives \(J_L\) a commutative group scheme
  structure over \(L\), and \(J_L\to\Spec L\) is of finite type.
\end{lemma}

\begin{proof}
  For each \(L\)-scheme \(T\), transport addition, zero, and inverse from
  \(\operatorname{Pic}^0_{C_L}(T)\) to \(\Hom_L(T,J_L)\) through
  \(r_L(T)\). Naturality of \(r_L\) and of pullback of Picard classes makes
  these operations natural in \(T\). The product of the represented functor
  with itself is represented by \(J_L\times_L J_L\), and the one-point
  functor is represented by \(\Spec L\). Yoneda therefore gives morphisms
  \(m:J_L\times_L J_L\to J_L\), \(e:\Spec L\to J_L\), and
  \(\iota:J_L\to J_L\). Each group identity holds on every functor of
  points, so faithfulness of Yoneda gives the corresponding equality of
  morphisms. Finally, Theorem~\ref{route-thm-glue-package} gives local
  finite type and quasi-compactness of this carrier. Together these are
  finite type.
\end{proof}

\begin{theorem}[Quasi-projectivity of algebraic groups: source theorem]
  \label{route-thm-algebraic-group-quasiprojective}
  \notready
  \source{stacks-project}
  \dcref{Tag 0BF7}

  If \(k\) is a field and \(G\) is a group scheme of finite type over
  \(k\), then \(G\) admits a locally closed immersion into
  \(\mathbf P^n_k\) for some finite \(n\).
\end{theorem}

This is the algebraic-group quasi-projectivity theorem of
\cite[Tag 0BF7]{stacks-project}, used here as a source theorem. Its
hypotheses do not require properness, smoothness, or geometric
connectedness. The source proof uses descent of the existence of an ample
invertible sheaf under field extension (Tag 0BDC), decomposition into
translates of the identity component (Tag 0B7R), smoothness in
characteristic zero and of the reduced subgroup over a perfect field
(Tags 047N, 047R, and 047P), and lifting the existence of an ample
invertible sheaf from the reduction of a Noetherian scheme in positive
characteristic (Tag 0BD4, case (3)). For the remaining smooth connected group
over an algebraically closed field, the source constructs an ample
invertible sheaf from translates of an effective Cartier divisor with
affine complement and constancy of the resulting Picard class in a
rational family (Tag 0BEH). These are the source prerequisites of this
theorem; no Picard representability theorem is one of them.

\begin{corollary}[Quasi-projectivity of the finite-stage carrier]
  \label{route-thm-exact-carrier-quasiprojective}
  \notready
  \uses{route-lem-exact-carrier-group,
    route-thm-algebraic-group-quasiprojective}

  The scheme \(J_L\) is quasi-projective over \(L\); thus, for some
  finite \(n\), there is a locally closed immersion
  \[
    j:J_L\hookrightarrow\mathbf P^n_L.
    \tag{QP}
  \]
\end{corollary}

\begin{proof}
  Apply Theorem~\ref{route-thm-algebraic-group-quasiprojective} to the
  finite-type group scheme of Lemma~\ref{route-lem-exact-carrier-group}.
\end{proof}

\begin{lemma}[Homogeneous prime avoidance]
  \label{route-lem-homogeneous-prime-avoidance}
  \notready
  \source{stacks-project}
  \dcref{Tag 00JS}

  Let \(S\) be a nonnegatively graded ring, let \(I\subseteq S_+\) be
  a homogeneous ideal, and let
  \(\mathfrak p_1,\ldots,\mathfrak p_r\), with \(r\geq1\), be
  homogeneous prime ideals. If \(I\nsubseteq\mathfrak p_i\) for every
  \(i\), there is a homogeneous \(f\in I\) of positive degree outside
  every \(\mathfrak p_i\).
\end{lemma}

\begin{proof}
  Discard primes contained in other primes in the list: avoiding the
  remaining primes avoids every discarded prime. They are now pairwise
  incomparable. For each \(i\), choose a homogeneous
  \(a_i\in I\setminus\mathfrak p_i\); its degree is positive because
  \(I\subseteq S_+\). For \(j\ne i\), choose a homogeneous
  \(b_{ij}\in\mathfrak p_j\setminus\mathfrak p_i\). Such homogeneous
  choices exist by taking a homogeneous component outside the indicated
  homogeneous prime. Put \(c_i=a_i\prod_{j\ne i}b_{ij}\). Then \(c_i\)
  is homogeneous of positive degree \(d_i\), lies in \(I\), is outside
  \(\mathfrak p_i\), and belongs to every other listed prime. Choose a
  positive common multiple \(d\) of the \(d_i\). The homogeneous element
  \[
    f=\sum_i c_i^{d/d_i}\in I
  \]
  has degree \(d\). Modulo \(\mathfrak p_i\) all terms except the
  \(i\)-th vanish, and that term is nonzero because \(\mathfrak p_i\)
  is prime. Thus \(f\) avoids every listed prime.
\end{proof}

\begin{lemma}[Finite subsets of quasi-projective schemes lie in affine opens]
  \label{route-lem-quasiprojective-finite-affine}
  \notready
  \source{stacks-project}
  \dcref{Tag 01ZY, case (4)}
  \uses{route-lem-homogeneous-prime-avoidance}

  Let \(X\) be quasi-projective over a field \(k\). Every finite subset
  of \(|X|\) is contained in an affine open of \(X\).
\end{lemma}

\begin{proof}
  Choose a locally closed immersion \(X\hookrightarrow\mathbf P^n_k\),
  put \(S=k[X_0,\ldots,X_n]\), and let
  \(Z=\overline X\setminus X\). This is closed in
  \(\operatorname{Proj}(S)\). Let \(E\subseteq|X|\) be finite. If
  \(E\) is empty the empty affine open suffices. Otherwise, for each
  point of \(E\), choose a standard open neighbourhood \(D_+(f_i)\)
  disjoint from \(Z\), with \(f_i\) homogeneous of positive degree.
  These finitely many opens cover \(E\). The ideal
  \(I=(f_1,\ldots,f_s)\subseteq S_+\) is not contained in any of the
  homogeneous primes corresponding to points of \(E\). By
  Lemma~\ref{route-lem-homogeneous-prime-avoidance}, there is a
  positive-degree homogeneous \(f\in I\) outside all these primes.
  Hence
  \[
    E\subseteq D_+(f)\subseteq\bigcup_i D_+(f_i)
      \subseteq\operatorname{Proj}(S)\setminus Z.
  \]
  The intersection \(X\cap D_+(f)\) is therefore a closed subscheme of
  the affine scheme \(D_+(f)\), and is the required affine open of \(X\).
  This argument uses homogeneous polynomials of unrestricted positive
  degree and applies over finite fields as well.
\end{proof}

\begin{theorem}[Affine neighbourhoods of finite-stage Galois orbits]
  \label{route-thm-exact-carrier-orbits}
  \notready
  \uses{route-thm-exact-carrier-quasiprojective,
    route-lem-quasiprojective-finite-affine, route-thm-finite-stage-yoneda}

  Let \(\rho\) be the semilinear action of
  \(\Gamma=\operatorname{Gal}(L/K)\) on \(J_L\) obtained from the
  canonical action on \(\operatorname{Pic}^0_{C_L}\) through \(r_L\).
  For every scheme point \(x\in|J_L|\), there is an affine open
  \(V\subseteq J_L\) such that
  \[
    \{\rho(\sigma)(x):\sigma\in\Gamma\}\subseteq V.
    \tag{OA}
  \]
\end{theorem}

\begin{proof}
  The group \(\Gamma\) is finite because \(L/K\) is finite Galois.
  Its orbit of \(x\) is thus a finite subset of the quasi-projective
  scheme of Corollary~\ref{route-thm-exact-carrier-quasiprojective}.
  Apply Lemma~\ref{route-lem-quasiprojective-finite-affine} to that subset.
\end{proof}

\section{The quotient and its represented functor}

\begin{theorem}[Conditional finite-Galois descent of the degree-zero functor]
  \label{route-thm-finite-galois-descent}
  \lean{AlgebraicGeometry.pic0SemilinearGalActionOfRepresentableBy,
    AlgebraicGeometry.pic0GaloisInvariantEquiv,
    AlgebraicGeometry.pic0GaloisInvariantEquivGaloisEquivariantOver,
    AlgebraicGeometry.pic0RepresentableBy_finiteGaloisDescent,
    AlgebraicGeometry.locallyOfFiniteType_pic0FiniteGaloisDescent,
    AlgebraicGeometry.quasiCompact_pic0FiniteGaloisDescent}
  \leanok
  \source{stacks-project}
  \dcref{Tags 0246 and 040L}

  Let \(L/K\) be finite Galois and let \(C/K\) be a smooth proper
  geometrically irreducible curve. Suppose that an \(L\)-scheme \(J_L\)
  represents \(\operatorname{Pic}^0_{C_L}\) through a natural
  equivalence \(r_L\), and that every orbit of the induced semilinear
  Galois action \(\rho\) on \(J_L\) lies in an affine open. Then the
  quotient obtained by descending Galois-stable affine opens is a
  \(K\)-scheme \(J\) representing \(\operatorname{Pic}^0_C\). More
  precisely, for every \(K\)-scheme \(T\) there are natural equivalences
  \[
    \Hom_K(T,J)
      \simeq\Hom_L(T_L,J_L)^{\operatorname{Gal}(L/K)}
      \simeq\operatorname{Pic}^0_{C_L}(T_L)^{\operatorname{Gal}(L/K)}
      \simeq\operatorname{Pic}^0_C(T).
    \tag{G}
  \]
  If \(J_L\to\Spec L\) is locally of finite type, respectively
  quasi-compact, then so is \(J\to\Spec K\).
\end{theorem}

\begin{proof}
  Yoneda transports the commutative group structure of the represented
  functor to \(J_L\); a group scheme over a field is separated. Given
  an affine open containing an orbit, intersect its finitely many Galois
  translates. This intersection is affine by separatedness, is
  Galois-stable, and still contains the orbit. These opens cover \(J_L\).
  On a stable affine open \(\operatorname{Spec}(B)\), finite-Galois
  descent of algebras identifies \(B\) with
  \(L\otimes_K B^{\operatorname{Gal}(L/K)}\). Descending the overlaps
  gives open immersions between the corresponding affine quotients;
  descent of morphisms preserves their cocycle equations. Glue along
  these immersions to obtain \(J\), with \(J\times_K L\cong J_L\).
  This is the stable-affine instance of effective fpqc descent in
  \cite[Tag 0246]{stacks-project}.

  Descent of morphisms, \cite[Tag 040L]{stacks-project}, gives the first
  equivalence in (G). The definition of \(\rho\) makes \(r_L\)
  equivariant, so restricting \(r_L(T_L)\) to invariant elements gives
  the second. The finite etale cover \(T_L\to T\) has double overlap
  the disjoint union of the Galois translates of \(T_L\). The sheaf
  condition for degree-zero Picard classes consequently identifies its
  matching families with the displayed invariant classes, giving the
  third equivalence. Smoothness of \(C\) implies the geometric
  reducedness required for this Picard comparison. All three
  equivalences commute with pullback along \(K\)-morphisms of tests,
  so their composite represents \(\operatorname{Pic}^0_C\).
  Local finite type and quasi-compactness descend along the faithfully
  flat quasi-compact cover \(\Spec L\to\Spec K\), giving the last
  assertions.
\end{proof}

\begin{corollary}[Original-field representability]
  \label{route-thm-original-field-representability}
  \notready
  \uses{route-thm-descended-universal-class,
    route-thm-finite-stage-yoneda, route-thm-exact-carrier-orbits,
    route-thm-finite-galois-descent, route-thm-glue-package}

  The stable-affine quotient \(J\) of the selected carrier \(J_L\)
  represents \(\operatorname{Pic}^0_C\) over \(K\) and is of finite
  type over \(K\).
\end{corollary}

\begin{proof}
  The joint-stage theorem supplies the finite Galois extension \(L/K\).
  Theorem~\ref{route-thm-finite-stage-yoneda} supplies \(r_L\) on
  \(J_L\), and Theorem~\ref{route-thm-exact-carrier-orbits} supplies
  its orbit condition. Apply
  Theorem~\ref{route-thm-finite-galois-descent} to these data.
  Local finite type and quasi-compactness of \(J_L\) follow from
  Theorem~\ref{route-thm-glue-package}; they descend to \(J\), proving
  finite type.
\end{proof}

\begin{remark}[Compatibility of the carrier throughout descent]
  \label{route-rem-no-carrier-switch}
  \uses{route-thm-original-field-representability}

  The class \(\xi_L\), its representing equivalence \(r_L\), the group
  structure, and the quasi-projective immersion all belong to the same
  \(L\)-scheme \(J_L\). The quotient is formed from the semilinear
  action defined by this \(r_L\). Quasi-projectivity supplies exactly
  the finite affine neighbourhoods needed for this descent; the argument
  makes no properness or projectivity assertion about \(J_L\).
\end{remark}
